\documentclass[a4paper,11pt]{article}

\usepackage{amsmath, amssymb, amsfonts}
\usepackage[margin=1cm]{geometry}
\usepackage{mathpazo}
\usepackage{eulervm}
\usepackage[T1]{fontenc}
\usepackage[utf8]{inputenc}
\pagenumbering{gobble}


\begin{document}

\begin{center}
  {\Large\textbf{Algebră și geometrie \\ Plan de învățămînt 2017 --- 2018 }}
\end{center}
\vspace{1cm}


\begin{itemize}
  \item Sisteme liniare și matrice (eliminare gaussiană, descompunere $A = LU, PA = LU$, cazul simetric $LDL^\prime$);
  \item Structuri algebrice, elemente de teoria grupurilor finite;
    \vspace{0.5cm}

  \item Spații și subspații vectoriale: independență liniară, generatori, bază, dimensiune;
  \item Transformări liniare și matrice asociate;
  \item Schimbarea bazei;
  \item Nucleul și imaginea;
  \item Exemple și aplicații;
    \vspace{0.5cm}

  \item Vectori și valori proprii pentru matrice pătrate;
  \item Cazul matricelor hermitice (simetrice);
  \item Diagonalizare prin matrice unitare (ortogonale);
  \item Diagonalizarea și jordanizarea matricelor pătrate complexe;
  \item Exponențiala unei matrice;
    \vspace{0.5cm}

  \item Geometrizarea spațiilor vectoriale: produs scalar și normă;
  \item Ortogonalizare Gram-Schmidt și descompunere QR;
  \item Aplicații;
    \vspace{0.5cm}

  \item Forme biliniare. Forme pătratice;
  \item Conice și cuadrice și reprezentare;
  \item Forme pătratice și matrice simetrice;
  \item Criterii de pozitiv definire (semidefinire);
    \vspace{0.5cm}
    
  \item Produse vectoriale și mixte;
  \item Planul și dreapta în spațiu;
    \vspace{0.5cm}

  \item Ecuația diferențială de ordinul 1. 
  \item Existența și unicitatea în condiții inițiale date;
  \item Proprietăți generale ale soluțiilor și exemple;
    \vspace{0.5cm}

  \item Ecuații diferențiale de ordinul $n$;
  \item Construcția soluției generale omogene și a unei soluții particulare omogene.
\end{itemize}

\newpage
\textbf{Bibliografie:}
\begin{itemize}
  \item Ana Niță și Alina Niță, \emph{Algebră liniară, Geometrie analitică}, Printech, 2009;
  \item Alina Petrescu-Niță, \emph{Matematici 2, Exerciții și probleme}, Printech, 2013;
  \item Octavian Stănășilă, \emph{Algebră liniară și geometrie}, vol. 1, ALL, 2000;
  \item Balea, S., Niță, Alina, Niță, Ana, Urseanu, R., \emph{Algebră liniară și geometrie}, Printech, 2004;
  \item Constantin, R., Moroianu, M., Mateescu, M., Niță, Ana, \emph{curs de algebră și geometrie}, Printech 2006.
\end{itemize}

\end{document}
